\documentclass[10pt]{article}

\usepackage[utf8]{inputenc}

\usepackage[T1]{fontenc}

\usepackage{amsmath}

\usepackage{amsfonts}

\usepackage{amssymb}

\usepackage[version=4]{mhchem}

\usepackage{stmaryrd}

\usepackage{mathrsfs}

\usepackage{graphicx}

\usepackage[export]{adjustbox}

\graphicspath{ {./images/} }



\begin{document}

\begin{enumerate}

  \setcounter{enumi}{3}

  \item С. с. - фундаментальный симплекс $\iota_{n}$ симплициального множества п. 3), к-рый в этом случае обозватается $\Delta_{n}$.

\end{enumerate}



СТАТИКА - раздел механики, в к-ром изучается равновесие материальных тел, находящихся под действием сил, и условия эквивалентности систем сил. Равновесие изучается по отношению к системе отсчета, в к-рой определены все силы, действующие на материальные тела (напр., равновесие по отношению к Земле). Кюак и в дина.ике, в С. вводятся модели реальных тел, дающие возможность сводить задачи о равновесии к более простым. Такими моделями являются материальная точка, упруго деформируемое тело, абсолютно твердое тело, континуум, идеальная жидкость, вязкая жидкость и т. П. В зависимости от свойств материальных тел С. разделяется на С. твердого тела, С. упруго деформируемого тела, С. Жидкости и газа. По своим методам исследования С. делится на геометрическую и аналитическую.



Ге о м е т р и ч е с к а я С. изучает геометрич. свойства систем сил (векторов сил), действующих на изучаемую систему материальных точек, построение эквивалентных систем сил, приведение их к простейшему виду и установление условий равновесия систем сил, а также тел, находящихся под действием сил. Важнейшей в геометрич. С. является С. абсолютно твердого тела.



В основе а в а л и т и ч е с к о й С. лежит понятие работы сил, действующих на систему материальных точек, на произвольном возможном перемещении системы. Основной принцип аналитич. С.- возможных перемешений приниип.



Начало развития С. относится к глубокой древности и связано с именем греч. механика и математика Архимеда. Основные законы равновесия геометрич. С. установлены нидерл. механиком С. Стевином (S. Stevin) и франц. механиком и геометром Л. Пуавсо (L. Poinsot). Первая общая формули ровка привципа возможных перемещений принадлежит И. Бернулли (J. Bernoulli). Ж. Лагранж (J. Lagrange) дал первое доказательство принципа возможных перемещений и получил следствия о равновесии системы.



Лum.: [1] П у а н с о Л., Начало статики, пер. с франц., М.-П., 1920; [2] Б е р н у л л и И., Избранные сочинения по механике, пер. с франц., М. - Л., 1937; [3] Л а г р а н Ж Ж.,

Аналитическая механика, т. 1, пер. с франц., М. Л., 1950; Аналитическая механика, т. 1 , шер. с франц., М.- Л., 1950 тельно моментов сил, в кн.: О с т р о г р а д с к и й М. В. Избранные труды, М., $1958 ;$ [5] ЖК у к о в с к и Й Н. Е., Теоретическая механика, 2 изд., М. - J., 1952; [6] Ч а П л Ыг и н С. А., Курсы лекций по теоретической механике, в кн.: Ги н С. А., Курсы лекций по теоретической механике , В Кн.:

ч а п л ы г и н С. А., Собр. соч., т. 4, М.-Ј., 1949; [7] А Пп е л ь П., Теоретическая механика, т. 1, пер. с франц., М., $1960 . \quad$ E. H. Березкин.



СТАТИСТКА - термин, употребляемый в математич. статистике для названия функций от результатов наблюдений.



Пусть случайвая величина $X$ принимает значения в выборочном пространстве $\left(\mathfrak{X}, \mathscr{B}, \mathrm{P}^{X}\right)$. Тогда любое $\mathscr{B}$-измеримое отображение $T(\cdot)$ пространства $\mathfrak{X}$ в некрое измеримое пространство $(\mathfrak{B}, \mathcal{A})$ наз. с т а т и с т ик о й, при этом распределение вероятностей статистики $T$ определяется формулой



\[

\begin{aligned}

\mathrm{P}^{T}\{B\}= & \mathrm{P}\{T(X) \in B\}=\mathrm{P}\left\{X \in T^{-1}(B)\right\}= \\

& =\mathrm{P} X\left\{T^{-1}(B)\right\}(\forall B \in \mathcal{A}) .

\end{aligned}

\]



П р и м е р ы. 1. Пусть $X_{1}, \ldots, X_{n}$ - независимые одинаково распределенные случайные величины, имеющие дисперсию. Тогда С.



\[

\bar{X}=\frac{1}{n} \sum_{i=1}^{n} X_{i} \quad \text { и } \quad s^{2}=\frac{1}{n-1} \sum_{i=1}^{n}\left(X_{i}-\bar{X}\right)^{2}

\]



суть несмещенные оценки для математич. ожидания $E X_{1}$ и дисперсии $D X_{1}$ соответственно. 2. Члены вариационного ряда



\[

X_{(1)} \leqslant X_{(2)} \leqslant \ldots \leqslant X_{(n)}

\]



построенного по наблюдениям $X_{1}, \ldots, X_{n}$ суть порлдковые статистики.



\begin{enumerate}

  \setcounter{enumi}{2}

  \item Пусть случайные величины $X_{1}, \ldots, X_{n}$ образуют стационарный случайный процесс, спектральная плотность к-рого есть $f(\cdot)$. В этом случае С.

\end{enumerate}



\[

I_{n}(\lambda)=\frac{1}{2 \pi n}\left|\sum_{k=1}^{n} X_{k} e^{-i k \lambda}\right|^{2}, \quad \lambda \in[-\pi ; \pi]

\]



называемая периодограммой, при определенных условиях регулярности на $f(\cdot)$ является асимптотически несмещенной оценкой для $f(\cdot)$, т. е.



\[

\lim _{n \rightarrow \infty} E I_{n}(\lambda)=f(\lambda), \quad \lambda \in[-\pi ; \pi] .

\]



В теории оценивания и проверки статистич. гипотез важную роль играет понятие достаточной статистики, к-рая позволяет производить редукцию данных без потери инфорормации об изучаемом параметрич. семействе распределений.



Лuт.: [1] J е м а н Ә., Проверка статистических гипотез, пер. с англ., М., 1964. П М. С. Никулин.



СТАТИСТИЧЕСКАЯ ГИПОТЕЗА - определенное предположение о свойствах распределения вероятностей, лежащего в основе наблюдаемых случайных явлений. Результаты наблюдений представляются обычно в виде реализации нек-рой совокупности случайных величин, конечной или бесконечной. При этом совместное распределение этих случайных величин известно не полностью, и С. г. предполагает принадлежность его к нек-рому определенному классу распределений. В такой ситуации ставится задача статистических гипотез проверки.



СТАТИСТИЧЕСКАЯ ИГРА - антагонистическая uгра, в к-рой игрок I интерпретируется как природа, игрок II - как статистик, стратегия игрока I - как случайный процесс, стратегия игрока II - как решающее правило (решающая функция), а фуункция выигрыша игрока I - как функция потерь («риск») статистика. Смешанной стратегией игрока I будет тогда вероятностная мера на множестве случайных процессов, смешанной стратегией игрока II - рандомизированное решающее правило, а функция выигрыша игрока I определяется как математич. ожидание функции потерь статистика. С. и. возникли из задач математич. статистики (вапр., задачи оденки параметров, испытания гипотез и т. п.), в к-рых принимается минимаксный критерий (см. Минимакса принцип). Систематич. описание задач математич. статистики (в т. ч. задач последовательного авализа) как С. и. впервые было проделано А. Вальдом (A. Wald), показавшим справедливость теоремы о минимаксе для широкого класса С. и. (см. [1]). Эта теорема дает метод решения задач математич. статистики, т. к. в ряде случаев нахождение максимина удобнее и легче, чем нахождение минимакса, что, в свою очередь, облегчает нахождение оптимального рандомизи рованного правила.



Jum.: [1] В а л ь д А., Статистические решающие функции, в сб.: Позиционные игры, М., 1967; [2] Б л е к у э л л Д.,



\begin{center}

\includegraphics[max width=\textwidth]{2023_07_09_e216966a503f0753d85dg-1}

\end{center}



СТАТИСТИЧЕСКАЯ ОЦЕНКА - функция от случайных величин, применяемая для оценки неизвестных параметров распределения вероятностей. См. Оценка статистическая.



СТАТИСТИЧЕСКАЯ СУММА - функция, используемая в равновесной статистич. физике, равная нормировочной константе в выражении для плотности (или матрицы плотности - в случае квантовой системы) в каноническом гиббсовском ансамбле.



\begin{enumerate}

  \item В случае классич. системы плотность распределения Гибббса $p(\omega), \omega \in \Omega(\Omega-$ фазовое пространство

\end{enumerate}



\end{document}